% latexmk -pdf -pvc blatt7.tex
\documentclass{hott-übung}

\begin{document}

\setcounter{blattnummer}{7}
\setcounter{aufgabennummer}{0}

\blatt

\aufgabe{}
Zeige, dass Komposition von punktierten Abbildungen assoziativ ist.
Gehe dazu wie folgt vor: Übersetze in abhängige Typen.

\aufgabe{Strukturidentitätsprinzip für Monoide}
  Seien \((M,\cdot,e)\) und \((N,\cdot,e)\) Monoide.
  Eine Abbildung \(f : M \to N\) heißt \emph{Homomorphismus}, falls sie ein Homomorphismus von Halbgruppen \((M,\cdot) \to (N,\cdot)\) ist und \(f(e)=e\) erfüllt.
  Ist \(f\) außerdem eine Äquivalenz, so nennen wir \(f\) einen \emph{Isomorphismus}.
  Zeige:
  \[((M,\cdot,e) =_{\mathrm{Mon}} (N,\cdot,e)) \simeq \Iso((M,\cdot,e),(N,\cdot,e))\]


\aufgabe{}
  Sei \(G\) eine Gruppe und \((BG,\ast)\) eine Entschleifung von \(G\).
  Es bezeichne \(G\text{-}\Set\) den Typ der \(G\)-Mengen.
  \begin{enumerate}
    \item 
  \end{enumerate}
Gruppenoperationen sind abhängige Typen über BG, sigma ist homotopiequotient und pi Fixmenge.

\end{document}
